\documentclass[11pt,a4paper]{article}
\usepackage[margin=24mm]{geometry}
\usepackage[T1]{fontenc}
\usepackage{lmodern}
\usepackage{amsmath,amssymb,booktabs,longtable,array,hyperref}
\hypersetup{hidelinks}
\newcommand{\La}{\textbf{(a)}\ }
\newcommand{\Lc}{\textbf{(c)}\ }
\newcommand{\Ld}{\textbf{(d)}\ }
\newcommand{\scr}{\texttt{audit\_\allowbreak defb\_\allowbreak boundary.py}}
\newcommand{\file}[1]{\texttt{\detokenize{#1}}}
\newcolumntype{L}[1]{>{\raggedright\arraybackslash}p{#1}}
\emergencystretch=3em
\title{Independent audit of the defect boundary package:\\exact traces, universal holonomies, and an unjustified bulk-scheme selection}
\author{Quantum Boussinesq solitons project}
\date{4 November 2026}
\begin{document}
\maketitle
\begin{abstract}
\Lc The bulk-counterterm triangle shift is correct. Its interpretation as a
physical selection of $\lambda_g=0$ is refuted: fixed classical absorption
rapidities are an additional assumption, and a consistent displacement of the
kinematics cancels the counterterm contribution. Moreover $\lambda_g=0$ does
not give zero vacuum sound shift in the controlling good-equation calculation.
The printed fixed-momentum transmission sign also contradicts its own
classical phase. The elementary boundary Casimir formulas, corrected energy-trace sign,
stationary-root envelope theorem and stable-root counting survive independent
computation. The purported quantum uniqueness of the finite boundary
counterterms does not follow from a classical scattering-length calculation.
A new exact logarithm identity proves the canonical energy holonomy is nonzero
on every nondegenerate real KP triangle; independent interval arithmetic
certifies the stated momentum zero. The absorbed-momentum proof needs a
paired-branch-point correction, and its extension to every positive
$\lambda_g$ is false. \Ld A physical renormalised quantum defect theory,
quantum absorption locus and fusion amplitude remain unconstructed.
\end{abstract}

\section{Scope, independence and verdict table}
Labels mean (a) established in the literature, (b) literature heuristic,
(c) project derivation/computational evidence, and (d) unestablished physical
interpretation or research judgment. ``Confirmed'' always refers to the
specified mathematical statement and domain. A conditional continuum trace
is not a constructed quantum ground-state charge. \Ld

\Lc This audit reads \file{AGENTS.md}, \file{README.md}, the template
\file{audit_phase2i.tex}, the boundary package source and PDF, both package
scripts, \file{audit_defq.tex}, relevant errata N07 and earlier qualifications,
\file{good_renorm.tex}, its audit, and the PI's re-ranked programme in
\file{report.tex}. The new checker imports no repository module, reads no
numerical cache and reuses no package computational routine. It derives
sewing and quadratic operators from the local action in project variables;
uses physical-$x$ absorption quadrature, Hessian polarization, KP cross-ratio
phase derivatives, the map $q=y/(1-y)$, independent Gauss--Legendre meshes,
and a rational half-angle interval partition. Symbolic checks use SymPy;
numerical calculations normally use 80 decimal digits and tolerance
$10^{-60}$; interval operations use 45 digits with interval-valued constants.
Explicit cutoff and printed-table tolerances are stated separately.

\Lc Only the three requested audit-prefixed repository files are created.
No existing file, README, cache or git state is changed. The instruction not
to change git state overrides the repository's routine pull/commit workflow.
The two original scripts are executed separately, with only their final JSON
writes redirected to audit-prefixed files in \file{/tmp/}; no mathematical
statement or assertion is changed. The full old script suite is not run.

\begin{longtable}{@{}L{6.0cm}L{2.0cm}L{7.3cm}@{}}
\toprule Claim in the package & Verdict & Independent evidence and limitation\\\midrule\endhead
A1. Good/CZ dictionary; scattering $s$, spectrum $\eta s$ & Confirmed (c) & Action conversion, four sewing residuals and decaying-mode determinant; \scr\ §§1--3.\\
A2. Four classical charges and ``off-shell'' conservation & Qualified & Exact flux identities (c), conditional on the bulk equations and sewing, not unrestricted off-shell conservation.\\
A3. Genuine absorption and all six defq corrections & Confirmed (c) & Physical-$x$ quadrature, Hessian traces, covariance, asymptotic limits, polynomial roots and quotient.\\
A4. Exact $\Delta P_D^{\rm Cas}$, $\mathcal E_D^{\rm sub}$, $\Delta E_D^{\rm ren}$ & Confirmed (c) & Independent quadrature and partial fractions, including both orientations and $|s|>2$. Physical vacuum interpretation remains conditional.\\
A5. Linear UV term and ``unique'' subtraction & Qualified & Divergent coefficient and equal-mode endpoint confirmed (c); no unique finite scheme or exhaustive boundary-operator classification.\\
A6. Stationary-root envelope and vacuum-charge shifts & Confirmed (c) & Leibniz differentiation with independently varied endpoints and parameters.\\
A7. Conditions (i)--(iii) physically fix four zero counterterms & Refuted (c) & Reference-vacuum constraint couples $\delta\rho_0$ to $\delta s$; uncomputed loop slope controls $\delta s$. A zero-shift convention can be defined.\\
A8. Bulk $u^2$ charge shifts and triangle cube sum & Confirmed (c) & Profile moment, exact stiffness-dependent soliton charges and cubic factorisation.\\
A9. Defect integrability selects $\lambda_g=0$; ``zero sound shift'' & Refuted (c) & Shifted-locus cancellation; physical-momentum conservation differs from fixed classical edges; $s_1(0)=-9/(4\pi)$.\\
A10. One-loop transmission formula and numerical table & Refuted sign (c) & Table reproduced, but fixed-momentum sign contradicts the displayed classical phase. Physical determinant remains (d).\\
A11. At most one stable KP root; closed-cycle process & Qualified & Root count and injectivity confirmed (c); closed walks are necessary within that graph, not proof of a physical fusion channel.\\
A12. Elementary $L_1$, canonical charges, triangle formulas, symmetry and endpoint values & Confirmed (c) & Independent KP trace, differentiation, rational log identities and midpoint/endpoint checks.\\
A13. All-triangle $\eta\mathcal C_E<0$ and unique zero of $\mathcal C_P$ & Confirmed (c) & New exact energy proof and independent interval certificate for momentum. Package proof coverage is qualified.\\
A14. $\lambda_g$ invariance of $\mathcal C_P$ and shift of $\mathcal C_E$ & Confirmed (c) & Telescoping momentum and cubic energy identities; exceptional midpoint $\lambda_*$ explicitly checked.\\
A15. Absorbed weights and non-identical vanishing of $\mathcal R_{\rm abs}$ at $\lambda_g=0$ & Qualified & Signs confirmed by a stronger analytic bound (c); disjoint-branch premise false, repaired by paired coefficients. Not pointwise nonvanishing or physical no-go.\\
A16. Same strict weight signs for every $\lambda_g\ge0$ & Refuted (c) & At $\vartheta=\pi/6$, $\lambda_g=2$ reverses the signs.\\
A17. Pole-free ``smooth'' derivative and vertex telescoping & Qualified & Effective-phase derivative and telescoping confirmed (c); logarithmic singularities remain, and effective coefficients are not intrinsic defect charges.\\
A18. Complete quantum boundary theory / exhaustive quantum counterterms & Qualified (d) & The package itself leaves non-CZ operators and quantum resonances open; its abstract claims more than the calculation establishes.\\\bottomrule
\end{longtable}

\section{Highest priority: what the bulk counterterm does and does not select}
\paragraph{The shift formula is right (A8; \scr\ §1).}
\Lc For the project soliton $u=k^2\operatorname{sech}^2(kx/2)/2$,
substituting $y=\tanh(kx/2)$ gives
\begin{equation}
 \int_{\mathbb R}u^2\,dx=\frac23k^3.
\end{equation}
For $d_i=q_i-q_{i+1}$, $\rho_i=\eta\operatorname{sgn}d_i$,
$k_i=|d_i|$, and $d_1+d_2+d_3=0$, factorisation gives
\begin{equation}
 \sum_i\rho_i\delta E_{m,i}
 =\frac{2\lambda_g}{3}\eta\sum_i d_i^3
 =2\lambda_g\eta\,d_1d_2d_3.
 \label{eq:cubeshift}
\end{equation}
It is nonzero for a nondegenerate triangle and $\lambda_g\ne0$.
There is no sign or factor error in this particular package formula.

\Lc A direct independent route to the canonical variations is to introduce
bulk stiffness $a=1+2\hbar\lambda_g$, keeping the normalization of inertia
and $u^3$ fixed. The exact classical soliton has $k^2=a-v^2$ and
\begin{equation}
 P_a=\frac23v(a-v^2)^{3/2},\qquad
 E_a=\frac2{15}(a-v^2)^{3/2}(a+4v^2).
 \label{eq:stiffness}
\end{equation}
Differentiating at $a=1$ gives the package's canonical shifts
\begin{equation}
 \delta P_c=2\lambda_gvk,\qquad
 \delta E_c=\frac23\lambda_g k(3-2k^2),\qquad
 \delta E_c-v\delta P_c=\frac23\lambda_g k^3.
 \label{eq:countercharges}
\end{equation}
The last expression is the energy shift at fixed physical momentum. It
survives changing between the mode and canonical charge descriptions.
Thus merely changing charge coordinates does \emph{not} remove the
counterterm's physical fixed-momentum dispersion change.

\paragraph{Why scheme selection does not follow (A9).}
\Lc Conservation concerns the physical momenta of an actual allowed
process. It does not require an unchanged set of three classical
absorption rapidities to remain an allowed quantum process. Even at
$\lambda_g=0$, the canonical triangle has nonzero energy holonomy below.
The two charge pairs obey $E_c-E_m=v(P_c-P_m)$, so shifting the classical
label to keep the same physical momentum changes the energy by the
corresponding classical derivative. A mode-edge identity cannot be
promoted to a Ward identity for physical fusion while discarding these
terms. \Ld Quantum sewing is needed to specify the allowed momenta.

\Lc There is a concrete kinematic counterexample to the inference from
\eqref{eq:cubeshift}. Scale an entire classical triangle as
$q_i(a)=\sqrt a\,q_i$, $k_i(a)=\sqrt a\,k_i$,
$v_i(a)=\sqrt a\,v_i$. Its charges scale as
$P_i(a)=a^2P_i(1)$, $E_i(a)=a^{5/2}E_i(1)$; their classical directed
sums stay zero. At first order the additional velocity shifts are
$\delta v_i=\lambda_gv_i$, and exact differentiation verifies
\begin{equation}
 \delta P_{c,i}+P_{0,i}'(v_i)\lambda_gv_i=4\lambda_gP_{0,i},\qquad
 \delta E_{c,i}+E_{0,i}'(v_i)\lambda_gv_i=5\lambda_gE_{0,i}.
 \label{eq:scaledclosure}
\end{equation}
Consequently the $\lambda_g$ contribution to the cycle cancels when its
kinematics are varied consistently. This checks the counterterm sector;
it does not erase the pre-existing one-loop canonical holonomy.
\Ld It is not a derivation of the actual quantum absorption shift or a
quantum-integrability proof for a particular boundary completion.

\Lc More generally, if a putative cycle has first-order remainders
$(B_E,B_P)$, two edge momentum displacements $x_1,x_2$ solve
\begin{equation}
 v_1x_1+v_2x_2=-B_E,\quad x_1+x_2=-B_P,
 \qquad
 x_1=\frac{-B_E+v_2B_P}{v_1-v_2},\quad
 x_2=\frac{B_E-v_1B_P}{v_1-v_2}.
\end{equation}
The first two triangle edges have distinct velocities and nonzero
classical momentum derivatives away from the degenerate endpoints.
Thus conservation alone cannot exclude a shifted locus.
\Ld Whether such shifts satisfy all quantum sewing and spectral conditions
is an additional, unsolved question.

\Lc Conversely, any single-valued boundary vacuum charge correction
$B(q)$ satisfies $\sum_i[B(q_{i+1})-B(q_i)]=0$. This holds outside the
CZ parameter family as well. Such a term cannot cancel the nonzero
fixed-edge sum \eqref{eq:cubeshift}. The same is true of a single-valued
shift of a vertex phase. \Ld More general boundary operators might change
the absorption locus, transmission, or admissible states; no computation
here proves or excludes that possibility. ``Outside CZ'' is therefore not
by itself a loophole in telescoping, nor is telescoping an exhaustive
quantum no-go theorem.

\paragraph{The good/bad and acoustic dictionaries (A1, A9).}
\Lc The retained good equation is
$u_{tt}=u_{xx}-u_{xxxx}-3(u^2)_{xx}$, with
$u=\varphi_x$ and $\Omega=\int d\pi\wedge d\varphi$.
The action substitution $U_{CZ}=-3\varphi$ gives $L_{CZ}=9L$ and
$D=D_{CZ}/9$. Scattering uses $s=\sigma_{\rm vac}=-2\sqrt3q_D$;
the decaying-mode spectrum uses $c=\eta s$. The package now implements
this distinction correctly, including absorption at both orientations.

\Lc The independently differentiated vacuum function
\begin{equation}
 W_g(U)=-\frac{(1-6U)^{3/2}-1+9U}{6\pi}
\end{equation}
gives $s_1=\lambda_g-9/(4\pi)$. Therefore the package's repeated name
``zero-sound-shift scheme $\lambda_g=0$'' is false. The conditional acoustic
choice in \file{good_renorm.tex} is $\lambda_g=9/(4\pi)$; this audit does
not adopt a scheme on the PI's behalf. The package's alternative description
of that value as ``Toda-matched'' is also inconsistent with the controlling
dictionary: $W_b(U)=-W_g(-U)$ implies $\lambda_b=-\lambda_g$, so
$\lambda_b=-3/(2\pi)$ continues to $\lambda_g=3/(2\pi)$.
The fixed spectral conventions further give $k_g=ik_b$,
$v_g(i\beta)=\epsilon\cosh\beta$ and
$\ell_b=-\ell_g(i\beta)$, not equality of the signed spectral parameters.
\Ld The conditional acoustic prescription remains a separate physically
motivated renormalisation choice; the defect calculation does not overturn it.

\section{Local charges, absorption and exact continuum traces}
\paragraph{Classical checks and the six corrections (A2--A3; \scr\ §§2--3).}
\Lc With $A=\eta/\sqrt3$, $\sigma=s+(b-a)/A$ and
$\rho(z)=\rho_0+(z-z^3/3)/(2A)$, our independent action uses
\begin{align}
 F&=A(a_x+2b_x)+\frac{\sigma^2}{12A},&
 G&=A(2a_x+b_x)+\frac{\sigma^2}{12A},\\
 E_D&=-A\sigma(a_x^2+a_xb_x+b_x^2)
           +\tfrac13\rho(\sigma)(a_x+b_x)-\tfrac19\tau(\sigma),\\
 P_D&=-\tfrac A2(a_x^2+4a_xb_x+b_x^2)+\tfrac19h(\sigma),&
 \Pi_D&=F+G,\quad M_D=b-a.
\end{align}
Here $\tau'=-z\rho$ and $h'=-\rho$. Variation of $D=a_tF+b_tG-E_D$
and the bulk action reproduces four sewing equations. Symbolic substitution
in the four independently differentiated flux balances gives zero for
arbitrary remaining jets, both orientations, and arbitrary integration
constants. This is conservation on the sewing/bulk equations, not conservation
for arbitrary off-shell fields.

\Lc Physical-$x$ integration, at $t=-2,0.7,3$, gives a maximum
absorption-charge discrepancy $7.91\times10^{-82}$. The input obeys
$s=\eta\sqrt3k-v$ and has $z_t=0$. The previous wrong input still gives
$\sqrt{21}-3$, not absorption. Hessians of the full charges, contracted with
complex modes obtained from the linear sewing equations, confirm
\begin{equation}
 \langle E_{D,2}\rangle_q+\langle E_{D,2}\rangle_{-q}
 =\langle K\rangle_q+\langle K\rangle_{-q}
 =\frac{\omega}{4\pi}\{\theta_D'(q)+\theta_D'(-q)\}.
 \label{eq:plus}
\end{equation}
The plus sign is correct. Plane-wave integration-by-parts boundary terms
vanish. The covariance is $1/(4\pi\omega)$; time averaging a real cosine
of amplitude $1/\sqrt{\pi\omega}$ instead gives twice this value.
Both signed-wavenumber transmission limits are $e^{-2\pi i\eta/3}$.
Fresh exponential-mode sewing determinants reproduce the corrected
$-2S^2(S^2-1)P_-P_+/3$ factor after matching row order. Direct roots at
$c=\sqrt3$ give one positive rate $\sqrt3/4$ and a marginal rate zero.
Finally $(\ell(p)-\ell(q))/(p-q)$ is the polynomial
$3\sqrt3[4(p^2+pq+q^2)-1]$, with no simple pole.

\paragraph{Elementary charges rederived (A4; \scr\ §4).}
\Lc Put $D_\pm=4u^2\pm2su+s^2-3$, $u=\sqrt{1+q^2}$.
The Hessian contractions first reproduce the spectral densities pointwise.
Their rational forms then admit the independent partial fractions
\begin{align}
 \frac{2[(s^2-2)u^2+(3-s^2)/2]}{D_+D_-}
 &=\frac{-su-1}{2D_+}+\frac{su-1}{2D_-},\\
 \frac{-4u^2+(s^2-3)^2}{4D_+D_-}
 &=\frac{(s^2-2)u/(4s)+(s^2-3)/8}{D_+}
  +\frac{(2-s^2)u/(4s)+(s^2-3)/8}{D_-}.
\end{align}
The second decomposition uses $s\ne0$; its original integral is regular
at $s=0$. Completing the squares
$D_\pm=4(u\pm s/4)^2+3(s^2-4)/4$ gives logarithms below $|s|=2$
and arctangents above it, with the continuous repeated-root limit at $|s|=2$.
Define
\begin{equation}
 L(s)=\log\left|\frac{s+1}{s-1}\right|,\qquad
 B(s)=\begin{cases}
 -\sqrt{4-s^2}\log\left|\dfrac{\sqrt3+\sqrt{4-s^2}}{\sqrt3-\sqrt{4-s^2}}\right|,&|s|<2,\\
 0,&|s|=2,\\
 2\sqrt{s^2-4}\arctan\!\dfrac{\sqrt{s^2-4}}{\sqrt3},&|s|>2.
 \end{cases}
\end{equation}
The resulting package formulas are confirmed:
\begin{align}
 \Delta P_D^{\rm Cas}&=\frac{\eta}{8\pi}\,[\sqrt3sL(s)+B(s)],\\
 \Delta E_D^{\rm ren}&=\frac{\eta}{16\pi}
 [-4\sqrt3s+\sqrt3(2-s^2)L(s)+sB(s)],\\
 \mathcal E_D^{\rm sub}&=\Delta E_D^{\rm ren}+\frac{\sqrt3\eta s}{4\pi}.
 \label{eq:closed}
\end{align}
The displayed formulas exclude $s=\pm1$. On unstable vacua they are formal
continuum traces, not constructed ground-state expectations. \Ld

\begin{center}\small
\begin{tabular}{rrrr}\toprule
$s$ & $\eta$ & $\Delta P_D^{\rm Cas}$ & $\Delta E_D^{\rm ren}$\\\midrule
$-1.2$ & $1$ & $-0.00720648262722269233$ & $0.24243415056776965958$\\
$-1.5$ & $1$ & $0.06055750652521987455$ & $0.29997546003179577445$\\
$-2$ & $1$ & $0.15142417490204896669$ & $0.35137653516192050810$\\
$-2.5$ & $1$ & $0.23117575764210336842$ & $0.36217135204446018208$\\
$1.7$ & $-1$ & $-0.09894765410120971040$ & $0.32603918892805694405$\\
$2.6$ & $-1$ & $-0.24638175981914822630$ & $0.35997176172617868413$\\\bottomrule
\end{tabular}
\end{center}
\Lc The independent $q=y/(1-y)$ quadratures agree with the elementary
expressions to $7.97\times10^{-82}$ over nine tests, including $s=0,0.5$
and both orientations. This is numerical agreement at working precision,
not an $82$-digit rigorous error bound.

\paragraph{UV subtraction and error accounting (A5).}
\Lc Set $C=\sqrt3\eta s/(4\pi)$ and $T=\sqrt{1+\Lambda^2}$.
The local subtraction is $Cq/\sqrt{1+q^2}$. Equal-mode counting also retains
\begin{equation}
 B_\Lambda=\frac{\Lambda T}{4\pi}
 [\theta_D^{\rm asympt}(\Lambda)-\theta_D^{\rm asympt}(-\Lambda)].
\end{equation}
Since $B_\Lambda=C\Lambda+O(\Lambda^{-1})$, the renormalised equal-mode
limit is $\mathcal E_D^{\rm sub}-C$, as in \eqref{eq:closed}, whereas
$E_D^{\rm mode}=-2C(T-1)+\Delta E_D^{\rm ren}+O(\Lambda^{-1})$.
Our checker evaluates this endpoint directly from the continuous rational
transmission phase at two cutoffs; the observed remainder decreases as
$1/\Lambda$. The coefficient of the linear divergent vacuum term is fixed.
Its $M_D$ part completes the open-half-line tadpole into the topological
charge, as the package states. \Ld This does not uniquely determine the
finite subtraction, the off-vacuum boundary action, or all quantum operators.

\Lc Mesh, arithmetic and cutoff errors are separated as follows. For the
most demanding mesh test $s=-1.2$, three panels in $y$ with $16,32,64$
Gauss--Legendre nodes per panel give maximum charge discrepancies
$1.04\times10^{-17}$, $1.44\times10^{-35}$,
$1.21\times10^{-64}$ respectively. Changing adaptive arithmetic from 80 to
60 digits changes the momentum by $2.10\times10^{-64}$. These are independent
sensitivity diagnostics; no extrapolation or fit is used for boundary
charges, so there is no fit-spread contribution. The original package uses
one precision and one compact quadrature, reporting agreement with another
closed representation rather than a total error budget. Its impressive
roundoff-level discrepancies are credible as cross-checks, but not physical
uncertainties or certified quadrature bounds.

\section{The envelope theorem survives; quantum uniqueness does not}
\paragraph{Independent variation (A6; \scr\ §1).}
\Lc At a root $\rho(\sigma_i)=0$, endpoint variations of
$\tau=\tau_0-\int_s^{\sigma_i}z\rho(z)\,dz$ and
$h=h_0-\int_s^{\sigma_i}\rho(z)\,dz$ vanish at first order when
$\rho(s)=0$. The remaining variations are precisely
\begin{align}
 \delta E_{D,i}&=-\delta\tau_0/9+
       \tfrac23\delta\rho_0(q_i^2-q_D^2),\\
 \delta P_{D,i}&=\delta h_0/9+
       \tfrac{2\sqrt3}{9}\delta\rho_0(q_i-q_D),\qquad
 \delta q_i=\frac{\eta\delta\rho_0}{3(1-12q_i^2)}.
\end{align}
The root-shift formula assumes a simple root; cusp roots need a different
expansion. These are valid variations of the classical CZ family. They
are not a calculation of the one-loop effective boundary potential.

\paragraph{The missing coupling and the missing loop calculation (A7).}
\Lc Preserving the reference vacuum at $\sigma=s+\hbar\delta s$ actually
requires
\begin{equation}
 0=\delta\rho_0+\frac{\sqrt3\eta}{2}(1-s^2)\delta s,
 \label{eq:constraint}
\end{equation}
not $\delta\rho_0=0$ independently of $\delta s$. The latter follows only
if $s$ has already been held fixed. The envelope theorem removes the
\emph{direct endpoint} contribution of $\delta s$; it does not remove the
$\delta\rho_0$ contribution induced by \eqref{eq:constraint}.

\Lc The classical positive-momentum slope really is
$a_0(s)=-2\sqrt3\eta/(s+1)$, away from its cusp, and
$a_0'(s)=2\sqrt3\eta/(s+1)^2\ne0$. Hence it is an invertible coordinate
on each admissible stable half-line. But write the physical slope as
$a_{\rm phys}=a_0(s)+\hbar[a_{\rm loop}(s)+a_0'(s)\delta s]+\cdots$.
Matching it to $a_0(s_{\rm ren})$ fixes, in this convention,
\begin{equation}
 \delta s=-\frac{a_{\rm loop}(s)}{a_0'(s)}.
\end{equation}
Nothing in the package computes $a_{\rm loop}$ or proves it vanishes.
Its script differentiates the \emph{classical} transmission factor at zero
momentum and then prints a quantum-uniqueness conclusion. Nor does
$z_{\rm ph}(0^+)=1$ for this classical family prove a quantum topological
non-renormalisation theorem. Boundary tadpoles and quantum stationarity
could likewise enter the physical version of \eqref{eq:constraint}.

\Ld A scattering length is a natural observable with which to define a
renormalised defect parameter. Calling the resulting parameter $s$ is a
reasonable renormalisation convention; it is not evidence that its bare-to-
renormalised shift is zero. Conditions (i)--(iii) can be made to define a
zero-shift classical parametrisation and chosen additive charge origins.
They do not derive unique physical zero counterterms without the missing
loop amplitude and boundary effective action. Equating this prescription
to a renormalised absorption root is another unproved identification.

\section{Transmission and the stable-vacuum graph}
\paragraph{What the transmission calculation tests (A10; \scr\ §8).}
\Lc Our soliton phase derivative comes directly from the KP cross ratio
with real roots $(v/\sqrt3\pm k)/2$ and phonon roots
$\operatorname{sgn}(q)\sqrt{1+q^2}/(2\sqrt3)\pm iq/2$.
The defect phase is a continuous argument of its rational transmission,
referenced to the same constant at both infinities. This avoids the
package's elementary dipole-phase decomposition. At $\alpha=0.35$ it gives
\begin{center}\small
\begin{tabular}{rrrr}\toprule
$q_D$ & $I_{SD}$ & $I_{SD}-\Delta x_D\Delta P_m$ & $I_{SD}-\Delta x_D\Delta P_c$\\\midrule
$0.50$ & $-1.33210991172351758$ & $-2.83443578117603636$ & $-3.78971817125468003$\\
$0.52$ & $-1.36337559539974186$ & $-2.67213057224522436$ & $-3.50432723991219632$\\\bottomrule
\end{tabular}
\end{center}
\Lc The table is confirmed as an evaluation of the printed formula, but
that formula has a convention inconsistency. Independent symbolic
 differentiation of the package's displayed effective classical phase gives
\begin{equation}
 \boxed{\partial_P\delta_{D,\rm cl}=-\Delta x_D,\qquad
 -\Delta P\,\partial_P\delta_{D,\rm cl}=+\Delta x_D\Delta P.}
 \label{eq:phasesign}
\end{equation}
The original classical note \file{defect_semiclassics.tex} explicitly uses
this minus relation, whereas the inherited one-loop note switches it.
At $\alpha=0.4$, $q_D=0.51$, the discrepancy with the asserted positive
relation is $0.6373325868$ in magnitude for $\epsilon=-1$ and
$2.7141893147$ for $\epsilon=+1$; both defect orientations are checked
at 80 digits. An exact polynomial reduction on $k^2+v^2=1$ confirms the
negative relation universally away from singularities.

\Lc Thus a fixed-momentum Taylor expansion of the displayed classical
phase contributes $+\Delta x_D\Delta P$, not the table's negative term.
Keeping the displayed $I_{SD}$ convention for this algebraic comparison,
the counterterm contribution to that Taylor term is
$+2\lambda_gvk\Delta x_D$ for canonical momentum, and zero for
fixed-profile mode momentum. Reversing the entire classical phase convention
would also reverse the later vertex derivative; both printed signs cannot
be retained simultaneously. \Ld A complete physical phase requires a
consistent determinant convention, the boundary-compatible domain and the
quantum boundary conditions. The table neither supplies these nor proves a
physical scheme invariance. This sign refutation does not select a new
physical phase normalization on the PI's behalf.

\paragraph{Stability and cycles (A11; \scr\ §3).}
\Lc For fixed $\eta$, the relevant stable half-line is
$\eta q\ge1/(2\sqrt3)$. With $q=\eta[1/(2\sqrt3)+t]$,
$12q^2-1=4\sqrt3t+12t^2>0$ for $t>0$. The cubic is strictly injective
there, including the endpoint by continuity. Thus at most one of its
three roots is stable for that orientation. Fresh sewing determinants
and direct quadratic-root filtering reproduce the growing counts and the
$c=\sqrt3$ switch. Marginality is distinguished from strict stability.

\Lc Within a fixed three-vertex absorption graph, any walk whose initial
and final defect vacua are stable must return to the same vertex. The
boundary charge then cancels. This is a necessary bookkeeping condition,
and the three-edge cycle is one example; it does not exclude backtracking
or repeated closed walks. \Ld It does not construct the composite
$D(q_{\rm st})+S_1+S_2\leftrightarrow D(q_{\rm st})+\bar S_3$, establish
its amplitude, or show that the intermediate resonances and all bulk
external solitons are admissible quantum states. ``Stable vacuum'' also
has not been shown to mean a global quantum ground state. The source's
closed-cycle conclusion is therefore qualified, not a physical process
proved by monotonicity.

\section{All-triangle holonomy: a stronger independent proof}
\paragraph{Canonical charges from an independent trace (A12; \scr\ §§5--6).}
\Lc Integrating the KP cross-ratio phase derivative gives the mode charges.
The energy density subtraction and cutoff endpoint are retained explicitly.
At $k=0.23,0.61,0.94$, $v\Delta P_m-\Delta E_m$ agrees with
\begin{equation}
 L_1(k)=\frac{3k}{2\pi}-\frac{\sqrt3}{8\pi}(3-4k^2)
       \log\left|\frac{2k+\sqrt3}{2k-\sqrt3}\right|.
\end{equation}
Differentiation with $k=\sqrt{1-v^2}$ then yields the package's
$\Delta P_c=\partial_vL_1$ and $\Delta E_c=v\partial_vL_1-L_1$.
After independently subtracting the leading analytic tails, the
$\Lambda=10^{12}$ trace discrepancies are at most $2.76\times10^{-38}$;
for $k=0.61$ the $10^8/10^{12}$ cutoff spread is
$2.76\times10^{-26}$. The tolerance is $10^{-32}$ at the finer cutoff,
not the $10^{-60}$ used for convergent no-cutoff integrals. No fitted
coefficient is fed back into these charges. Their physical determinant
identification retains the earlier audits' qualifications. \Ld

\paragraph{An exact logarithm factorisation (A13).}
\Lc Set $t=\tan(\vartheta/2)$, $0<\vartheta<\pi/3$, and
$k_1=\sin\vartheta$, $k_2=\sin(\pi/3-\vartheta)$,
$k_3=\sin(\pi/3+\vartheta)=k_1+k_2$. Write $L_i$ for the three positive
logarithms in the package's holonomy formulas. Rational simplification gives
\begin{equation}
 e^{L_1}=\frac{(t+1/\sqrt3)(t+\sqrt3)}{(t-\sqrt3)(t-1/\sqrt3)},\qquad
 e^{L_2}=\frac{\sqrt3-t}{t(1+\sqrt3t)},\qquad
 \boxed{e^{L_3}=e^{L_1}e^{L_2},\quad L_3=L_1+L_2.}
 \label{eq:logs}
\end{equation}
All arguments are positive on the stated domain, so no $2\pi i$ ambiguity
enters this real-log identity. It reduces the canonical energy holonomy to
\begin{equation}
 \boxed{\quad
 \eta\mathcal C_E=-\frac{\sqrt3}{2\pi}
 [(k_3^2-k_1^2)L_1+(k_3^2-k_2^2)L_2]<0.
 \quad} \label{eq:strictenergy}
\end{equation}
Both differences are strictly positive since $k_1,k_2>0$ and
$k_3=k_1+k_2$. This is an all-triangle symbolic proof, not a sample or an
endpoint extrapolation. It alone proves the nonzero vector at $\lambda_g=0$.

\Lc The momentum expression reduces to
\begin{equation}
 \eta\mathcal C_P=\frac{\sqrt3}{\pi}
 [\sqrt3k_1L_2-(v_1-v_3)L_1]. \label{eq:CP}
\end{equation}
The exact rational reflection
$t\mapsto(1/\sqrt3-t)/(1+t/\sqrt3)$ exchanges $k_1,k_2$, negates the
appropriate velocities and proves $\mathcal C_E$ even and $\mathcal C_P$
odd about $\pi/6$. At the midpoint,
\begin{equation}
 \eta\mathcal C_E=-\frac{3\sqrt3}{4\pi}\log(2+\sqrt3),\quad
 \mathcal C_P=0,\quad
 \eta\mathcal C_P'=\frac{3\sqrt3\log(2+\sqrt3)-9}{\pi}.
\end{equation}
The static third edge there requires the continuous $v=0$, $k=1$
extension of the introductory open rapidity domain; the formulas are
regular at that point.

\Lc To certify the absence of another momentum zero, our checker uses
\eqref{eq:CP} in $t$, not the package's uniform $\vartheta$ partition.
On $0<t\le10^{-4}$, the mean-value bound for $L_1/t$, computed from the
exact derivative of its rational argument, gives
$\eta\mathcal C_P/t>14.8187$.
On $[2-\sqrt3-0.0005,2-\sqrt3]$, interval differentiation gives
$d(\eta\mathcal C_P)/dt<-1.2512$.
The remaining compact interval is covered by 47 adaptive dyadic cells;
every lower bound is positive, the smallest exceeding
$6.89\times10^{-5}$. Exact reflection covers the other half. Hence the
unique zero is $\vartheta=\pi/6$. This is a computer-assisted interval
proof with analytic endpoint control, not a wholly symbolic proof.

\Lc The package does more than sample: its interior interval tests are
real enclosures. But its tiny-endpoint section inserts, rather than derives
in executable code, the interval \file{["2.309", "2.310"]} for a divided
logarithm and several coefficient bounds. Its asymptotic checks are numerical
at three tiny points. These omissions warrant qualification of its proof
coverage; they do not refute the signs, which the independent certificate
and \eqref{eq:strictenergy} now establish. The reported endpoint expansions
and midpoint values are also independently confirmed.

\paragraph{Counterterm dependence (A14).}
\Lc Telescoping gives $\sum_i\rho_iv_ik_i=0$ and therefore
\begin{equation}
 \mathcal C_P^{\lambda_g}=\mathcal C_P,\qquad
 \mathcal C_E^{\lambda_g}=\mathcal C_E+\eta\lambda_g\sin3\vartheta.
\end{equation}
These are correct. At $\vartheta=\pi/6$, the exceptional value
$\lambda_*={3\sqrt3}\log(2+\sqrt3)/(4\pi)$ cancels the entire canonical
vector. It differs from both $0$ and $9/(4\pi)$, as the package notes.
The zero-scheme all-triangle theorem should not be extended to arbitrary
$\lambda_g$ without this exception.

\section{Absorbed momentum and the derivative's remaining singularity}
\paragraph{Weights: a stronger bound and a false extension (A15--A16).}
\Lc For $0<\alpha<\pi/3$, put $k=\sin\alpha$, $c=\cos\alpha$ and
$L=\log[(\sqrt3+2k)/(\sqrt3-2k)]$. Independent reduction gives
\begin{equation}
 w_m=\frac{\Delta P_m(\alpha,+1)}{k}
 =-\frac{\sqrt3c}{4\pi k}L
  +\frac1{4\pi}\log\frac{2c+1}{2c-1}
  +\frac1\pi\log\tan\frac\alpha2.
\end{equation}
Its derivative simplifies exactly to
\begin{equation}
 w_m'(\alpha)=\frac{\sqrt3 L}{4\pi\sin^2\alpha}>0,\qquad
 \lim_{\alpha\uparrow\pi/3}w_m(\alpha)=-\frac{3\log3}{4\pi}.
\end{equation}
This proves a stronger bound than the printed $-\log3/(2\pi)$.
The canonical weight is manifestly negative for positive velocity;
orientation reversal negates both momenta. Thus $w_1<0<w_2$ on every
triangle in the zero-counterterm scheme is confirmed analytically.

\Lc The package's upper-half interval function is not exactly $w_m$:
its final term has $k\log\tan(\alpha/2)/(2\pi)$ instead of
$k\log\tan(\alpha/2)/\pi$ before division by $k$. Its computed quantity
exceeds the true weight by $-\log\tan(\alpha/2)/(2\pi)>0$.
Therefore its negative upper bound can still imply negativity of the true
weight, but the claimed equality is false. Its 500 cells omit both open
endpoint strips, and its printed exact bound is not the same as its
asserted $-0.17$ bound. The new monotonicity proof resolves these gaps.

\Lc The extension to every $\lambda_g\ge0$ is refuted. Canonical weights
change by $2\lambda_gv_i$. At the symmetric triangle and $\lambda_g=2$,
\begin{equation}
 w_1^{\lambda_g}=2.20649946030080639>0,\qquad
 w_2^{\lambda_g}=-w_1^{\lambda_g}<0.
\end{equation}
At an intermediate positive value the two weights coincide. The package's
sentence that the added term is negative only if $\lambda_g<0$ also ignores
the negative velocity of the second edge. Its script contains no assertion
testing this extension. The mode weight, defined at fixed profile, is
unchanged by this particular bulk counterterm; that does not rescue the
claimed statement for both charge choices.

\paragraph{Branch points occur in pairs (A15).}
\Lc The claim that the logarithm associated with $q_2$ has branch points
disjoint from $q_1,q_3$ is false on the test-soliton root locus. If
$p_a=q_i$, the condition $p_a^2+p_ap_b+p_b^2=1/4$ forces $p_b$ to be
one of the other two triangle roots. Two vertex logarithms then become
singular simultaneously, with opposite logarithmic coefficients.
For $a_i=w_{i-1}-w_i$, their coefficients are $a_i-a_j$, not a lone $a_i$.

\Lc This repairs rather than destroys the functional obstruction:
all three paired coefficients vanish only if $a_1=a_2=a_3$; since
$\sum_i a_i=0$, this requires $w_1=w_2=w_3$, impossible when
$w_1<0<w_2$. At least one permitted signed test-soliton branch therefore
has a nonzero logarithmic coefficient on every nondegenerate triangle.
If the third edge is static at the midpoint, the other two already have
nonzero coefficients. Thus the remainder is not identically zero as a
function on the signed test-soliton domain at $\lambda_g=0$.
\Ld This is not a statement that its value is nonzero at every test
soliton, nor that a physical quantum fusion amplitude must fail. No such
pointwise or physical conclusion follows from non-identical vanishing.

\paragraph{Pole-free is not smooth (A17; \scr\ §7).}
\Lc Differentiating the effective classical phase does give
\begin{equation}
 \partial_q\delta_{D,\rm cl}
 =-\frac{2\sqrt3}{3}(12q^2-1)\epsilon\log|z_t|
       +4\epsilon\eta k(v+2\sqrt3q).
\end{equation}
The polynomial quotient cancels the simple poles. However the logarithm
still diverges at a generic absorption root, where $12q^2-1\ne0$.
The package's description as an ``exact smooth identity'' is therefore
false at those roots. It is a correct identity away from them with
logarithmic boundary behaviour.

\Lc The phase's $M_D=\epsilon\eta q$ and
$\Pi_D=\epsilon\eta(6q^2-1)/\sqrt3$ are the \emph{effective coefficients}
of \file{defect_semiclassics.tex}. They must not be confused with the
intrinsic vacuum charges from the local action,
$M_D^{\rm vac}=2\eta(q_D-q)$ and $\Pi_D^{\rm vac}=2\sqrt3\eta q^2$.
Our derivative check explicitly uses the effective phase definition.
The package's symbol reuse obscures that distinction. Its telescoping
identity for any single-valued vertex-phase shift is nevertheless exact.
\Ld It does not constrain all changes of quantum edge kinematics.

\section{Coverage, circularity, runtimes and failures recorded}
\paragraph{Which (c) statements do the original scripts actually check?}
\Lc The classical flux/absorption assertions, trace sign, envelope
polynomials, closed-form quadratures and interior interval inequalities are
substantive tests. The holonomy vertex-sum comparison substitutes the same
canonical charge formulas on both sides; it checks assembly and conventions,
not an independent derivation of their quantum physical meaning. Our KP
trace and canonical differentiation supply an independent route.
The bulk cube identity is tested, but the subsequent physical-selection
sentence is a printed inference. The slope calculation is classical;
no loop correction or quantum uniqueness system is solved. The transmission
section prints \file{[PASS]} after evaluating numbers without a comparison
assertion. The interval code never tests the positive-$\lambda_g$ extension,
the disjoint-branch-point premise, smoothness at absorption, or compatibility
of the fixed-momentum sign with its own classical phase.

\Lc The package does repair the six concrete defq errors. It does not
repair the earlier conditional identification of a physical charge pair,
the incomplete determinant/domain problem, or the absence of a quantum
sewing computation. These are controlling qualifications in the earlier
audits and errata, not conclusions superseded by another passing script.
No evidence from this snapshot establishes hidden historical tolerance
changes. The package records prior development failures honestly, but its
historical wrong-input comparison now uses $10^{-45}$ whereas its narrative
mentions an earlier $10^{-55}$ test. The new independent test uses
$10^{-60}$ and passes; the weaker current package threshold is disclosed
rather than used as evidence of fraud or a numerical failure.

\paragraph{Runtime and reproducibility.}
\Lc The initial isolated runs of the unchanged originals both exit zero
and print \file{ALL CHECKS PASSED}: renormalisation 27.91 seconds and
holonomy 9.59 seconds, run concurrently. The final independent checker has no cache and completes in
78.62 seconds. A separate expanded run including sequential original-script
checks measured 23.32 seconds for renormalisation and 8.26 seconds for
holonomy. The final checker reports its own runtime and reruns the originals
only with \file{--package-runs}. All independent assertions, including assertions that refuted
claims have a nonzero discrepancy, precede its final line
\file{ALL CHECKS PASSED}. No passing marker is manufactured after an error.

\paragraph{Failures recorded.}
\begin{enumerate}
\item \Lc Package refutations: physical selection of $\lambda_g=0$;
its zero-sound-shift and Toda-matching names; independent $\delta\rho_0=0$
from condition (i); quantum zero parameter shift inferred from a classical
slope; positive-$\lambda_g$ weight-sign extension; disjoint branch points;
smoothness at absorption; and the fixed-momentum transmission sign relative
to the displayed classical phase. The upper-half interval formula also has a
missing factor of two. These discrepancies are retained and, where
algebraic, asserted in the independent script.
\item \Lc Package proof/coverage gaps: uncomputed quantum slope and
boundary effective action, endpoint interval constants supplied rather than
derived, absent endpoint coverage in the mode-weight scan, numerical table
prints without assertions, and effective versus intrinsic defect-charge
notation. The independent exact logarithm and weight-monotonicity proofs
repair the mathematical sign results without proving their physical
interpretations.
\item \Lc The first independent run failed the determinant-prefactor
comparison by a sign ($+0.00174579217008414839$ versus its negative).
The audit ordered rows $(a_{xx},b_{xx},a_{xxx},b_{xxx})$, whereas the
quoted determinant uses the two left rows followed by the two right rows.
Swapping rows two and three to match that convention removed the discrepancy.
The target prefactor, polynomial and tolerance were unchanged. This was an
audit row-order error, not a refutation of the package determinant.
\item \Lc A late independent phase-derivative assertion against the
package's positive-delay relation failed by $0.63733258678449138541$.
The displayed classical phase instead obeys the negative-delay relation,
confirmed both symbolically and for all four direction/orientation choices.
The final checker asserts this nonzero discrepancy and the consistent
negative relation at the original $10^{-60}$ tolerance. No tolerance was
loosened and the failed equality is recorded as a package refutation.
\item \Lc Two exploratory repository searches named a nonexistent old
script/glob and returned a missing-path diagnostic. The relevant definitions
were read from the actual TeX sources. No numerical claim depended on those
failed searches. The first PDF build also failed because the default font setup requested
an unavailable bitmap font and its font cache was read-only. Selecting
the installed Latin Modern outline fonts fixed the build without altering
mathematics. The successful high-precision run retained all mathematical
assertions; none was disabled or weakened to obtain a pass.
\end{enumerate}

\paragraph{References and reading scope.}
\La Classical source: Corrigan--Zambon,
\emph{Adding integrable defects to the Boussinesq equation},
\href{https://arxiv.org/abs/2306.03772}{arXiv:2306.03772}.
The entry, abstract and full-text section 2 were read online; the entire
article was not read. The package's title differs from the verified entry.
No new Faddeev--Korepin reading or attribution is claimed.
\Lc New results are checked as project statements. The PDF was compiled
with two \file{pdflatex} passes and all pages rendered and inspected.
\end{document}
